\documentclass[main]{subfiles}
\usepackage[]{workbook}
\numbering{ws2-12}

\begin{document}

\ifSubfilesClassLoaded{%
  \chaptermark{\chaptertwo}
}{}

\section{Power Series Representations of Functions} \label{ws2-12}
\begin{framed}
\centerline{\textbf{Lesson Goals}}

You should be able to:
\begin{itemize}
\item Manipulate a known power series representation using legal algebraic operations to find other power series representations, and determine how each manipulation affects the radius and interval of convergence
\item Describe the relationship between a Taylor series and a power series representation, and use it to find the Taylor series of other functions more easily
\end{itemize}
\end{framed}



\begin{enumerate}

\item 
\begin{problem}[2.5in]
Use the Taylor series for $\frac{1}{1-x}$ to find a power series representation of $\displaystyle \frac{x^3}{1+27
x^3}$.  For what $x$ can we be sure that this representation is
valid?
\end{problem}

\begin{solution}
Since we are asked for a power series representation, we ought to
start with a power series representation we know and manipulate it.
In this case, let's start with the power series representation of
$\frac{1}{1 - u}$ that we know:
\[ \frac{1}{1 - u} = \sum_{k=0}^\infty u^k = 1 + u + u^2 + u^3 +
\cdots \textrm{ for } |u| < 1.\] If we substitute $u = -27 x^3$,
then the previous equation becomes
\[ \frac{1}{1 + 27x^3} = \sum_{k=0}^\infty (-27x^3)^k = 1 + (-27x^3) +
(-27x^3)^2 + \cdots \textrm{ when } |-27x^3| < 1.\] Saying $|-27 x^3|
< 1$ is the same as saying $\left| x^3 \right| < \frac{1}{27}$, or
$|x| < \frac{1}{3}$.  It's helpful to simplify the previous equation
before we keep going:
\[ \frac{1}{1 + 27x^3} = \sum_{k=0}^\infty (-27)^k x^{3k} = 1 - 27x^3
+ 27^2 x^6 + \cdots \textrm{ when } |x| < \frac{1}{3}.\] Now, let's
multiply both sides by $x^3$ (which does not change where the
representation is valid):
\[ \boxed{\frac{x^3}{1 + 27x^3} = \sum_{k=0}^\infty (-27)^k x^{3k + 3} =
x^3 - 27x^6 + 27^2 x^9 + \cdots \textrm{ when } |x| < \frac{1}{3}}\]
\end{solution}


\iffalse
\item 
\note{I don't expect to get to this in class, but I've put it here just
so they have an example where the radius of convergence is 0.}

In this problem, you'll look at the power series $\displaystyle
\sum_{n=0}^\infty n! x^n$.

\begin{enumerate}
  \item
  \begin{problem}[0.35in]
  Write out the first few terms of the power series.  Does the series
  converge when $x = 0$?
  \end{problem}

  \begin{solution}
  The power series $\displaystyle \sum_{n=0}^\infty n! x^n$ is equal
  to $0! + 1! x + 2! x^2 + 3! x^3 + \cdots$, or $\boxed{1 + x + 2! x^2
  + 3!  x^3 + \cdots}$.  When $x = 0$, this is equal to $1 + 0 + 0 +
  \cdots$, which \fbox{converges}.
  \end{solution}


  \item
  \begin{problem}[0.5in]
  Find all values of $x$ for which the series $\displaystyle
  \sum_{n=0}^\infty n! x^n$ converges.  (Make sure your answer is
  consistent with your answer to
  {\#8}.)
  \end{problem}

  \begin{solution}
  Let's try the Ratio Test: $\displaystyle \lim_{n \to \infty} \left|
  \frac{(n + 1)! x^{n+1}}{n! x^n} \right| = |x| \lim_{n \to \infty}
  (n + 1)$.  If $x \neq 0$, then this limit is $\infty$, so the Ratio
  Test says that the series diverges.  If $x = 0$, then the limit is
  just $\displaystyle \lim_{n \to \infty} 0 = 0$, and the Ratio Test
  says that the series converges.  (This agrees with our answer to
  {{\#8}(a)}.)

  So, the series converges \fbox{only for $x = 0$}.
  \end{solution}

  \item
  \begin{problem}
  What is the radius of convergence of this power series?      
  \end{problem}

  \begin{solution}
  It is \fbox{0}, since the power series converges only at the center
  (that is, at $x = 0$).
  \end{solution}

  
\end{enumerate}
  \fi





 \begin{problemonly} 
  \begin{framed}
  \textbf{Theorem.} If the power series, \[\displaystyle \sum_{k=0}^\infty
  a_k (x - c)^k = a_0 + a_1 (x - c) + a_2 (x - c)^2 + \cdots\] has radius
  of convergence $R$, where $R > 0$ or $R = \infty$, then as a function
  $f(x) = \displaystyle \sum_{k=0}^\infty a_k (x - c)^k$ is differentiable on $(c - R, c + R)$.
  On that interval,
  \begin{enumerate}[topsep=0pt,itemsep=0pt]
  \item $\displaystyle f'(x) = \sum_{k=1}^\infty k a_k (x - c)^{k-1} =
  a_1 + 2 a_2 (x - c) + 3 a_3 (x - c)^2 + \cdots$. 
  \item $\displaystyle \int f(x)~dx = C + \sum_{k=0}^\infty
  \frac{a_k}{k + 1} (x - c)^{k + 1} = C + a_0 (x - c) + \frac{a_1}{2}
  (x - c)^2 + \frac{a_2}{3} (x - c)^3 + \cdots$.
\end{enumerate}
The power series in (a) and (b) both have radius of convergence $R$.
However, the interval of convergence may change.
\end{framed}
\end{problemonly}
\pspace{\newpage}

  \item 
  \begin{problem}[1.15in]
  Find a power series representation of $\displaystyle \frac{1}{(1 -
  x)^2}$ centered at $0$. What is the
  radius of convergence of the power series you have found?
  
  Hint: $\displaystyle \frac{1}{(1 - x)^2}$
  is the derivative of $\displaystyle \frac{1}{1 - x}$.
  \end{problem}

  \begin{solution}
  We know that
  \[ \frac{1}{1 - x} = 1 + x + x^2 + x^3 + \cdots = \sum_{k=0}^\infty
  x^k \textrm{ for } |x| < 1,\] and the radius of convergence of the
  power series is 1.  Let's differentiate both sides; the theorem says
  that we differentiate the power series exactly the way we expect, so
  we get: 
  \[ \frac{1}{(1 - x)^2} = 1 + 2x + 3x^2 + \cdots = \sum_{k=1}^\infty k
  x^{k-1} \textrm{ on the interval } (-1, 1),\] and the power series
  $\displaystyle \sum_{k=1}^\infty k x^{k-1}$ has radius of convergence
  1.

  {\LARGE\Pointinghand} A few notes:
  \begin{itemize}
    \item When writing these equations, it's important to say where the
    sigma notation starts (so, in the first equation, we said the sigma
    notation started at $k = 0$).
    \item When we differentiate in sigma notation, we just differentiate
    the general term; in this case, the derivative of $x^k$ was
    $kx^{k-1}$.  However, we do need to be careful about where the sigma
    notation starts; in this case, the $k = 0$ term in the original
    series was $x^0 = 1$.  Since the derivative of $x^0$ is just 0, we
    leave it out of the sigma notation when writing the derivative;
    that's why the sigma notation for the derivative starts at $k = 1$. 
    (We can also just check with the $\cdots$ notation: the first term
    in the power series representation of the derivative is $1$, which
    is $kx^{k-1}$ when $k = 1$.)      
  \end{itemize}
  \end{solution}
  
\pspace{\newpage}

    \note{Make sure to emphasize that, when we integrate, there is a $+C$
    and that we need to solve for that $C$ to get a valid representation
    of $\ln(1 + x)$.}

\item 
  \begin{enumerate}
    \item
    \note{Make sure to emphasize that, when we integrate, there is a $+C$
    and that we need to solve for that $C$ to get a valid representation
    of $\ln(1 + x)$.}

    \begin{problem}[2.5in]
    Find a power series representation for $\ln(1 + x)$ centered at $0$.
    (Hint: What is the derivative of $\ln(1 + x)$?)
    \end{problem}

    \begin{solution}
    The derivative of $\ln(1 + x)$ is $\frac{1}{1 + x}$, and we can find
    a power series representation of $\frac{1}{1 + x}$ pretty easily.
    We start with
    \[ \frac{1}{1 - x} = \sum_{k=0}^\infty x^k = 1 + x + x^2 + \cdots
    \textrm{ when } |x| < 1.\]
    Replacing $x$ by $-x$ everywhere,
    \[ \frac{1}{1 + x} = \sum_{k=0}^\infty (-x)^k = 1 - x + x^2 - \cdots
    \textrm{ when } |-x| < 1.\]
    Simplifying,
    \[ \frac{1}{1 + x} = \sum_{k=0}^\infty (-1)^k x^k = 1 - x + x^2 -
    \cdots \textrm{ when } |x| < 1.\] The radius of convergence of this
    power series is still 1 (it's a geometric series with common ratio
    $-x$), so the theorem says we can integrate term-by-term on the
    interval $(-1, 1)$:
    \[ \ln (1 + x) = C + \sum_{k=0}^\infty (-1)^k \frac{x^{k + 1}}{k +
    1} = C + x - \frac{x^2}{2} + \frac{x^3}{3} - \cdots \textrm{ for }
    x \textrm{ in } (-1, 1).\]
    Moreover, the theorem tells us that the radius of convergence of
    this new power series is 1.  To find $C$, we'll plug in $x = 0$ (in
    general, plugging in the center of the series is the fastest way to
    find the constant of integration):
    \[ \ln(1) = C + \sum_{k=0}^\infty 0 = C.\]
    So, $C = 0$, and 
    \[ \ln(1 + x) = \sum_{k=0}^\infty (-1)^k \frac{x^{k + 1}}{k + 1} = x
    - \frac{x^2}{2} + \frac{x^3}{3} - \cdots \textrm{ for } x \textrm{
    in } (-1, 1).\]

    {\LARGE\Pointinghand} \textbf{Tip:} Before you integrate a power
    series, simplify it!  We can integrate $x^{\textrm{any power}}$
    easily, but if we want to integrate $(\textrm{something involving
    $x$})^{\textrm{a power}}$ without first simplifying, we need to
    take into account the Chain Rule.
    \end{solution}

    \item
    \begin{problem}[2in]
    What is the radius of convergence for the power series you have
    found?  What is the interval of convergence?
    \end{problem}

    \begin{solution}
    As we already stated, the radius of convergence is \fbox{1}.

    Since we know the radius of convergence is 1, we know that the
    series converges when $|x| < 1$ and diverges when $|x| > 1$.  We
    only need to see what happens when $x = \pm 1$.  In fact, we
    already did this during the lesson on the Ratio Test, where we
    found that the interval of convergence is \fbox{$(-1, 1]$}.

    As a result, the equation
    \[ \ln(1 + x) = \sum_{k=0}^\infty (-1)^k \frac{x^{k + 1}}{k + 1} = x
    - \frac{x^2}{2} + \frac{x^3}{3} - \cdots\] that we found in
    {{\#4}(a)} is also valid for $x = 1$.  (This is true
    by a different theorem, called the Abel Limit Theorem.)  
    \end{solution}

    \item 
    \begin{problem}[\vfill]
    We showed in class that the alternating
    harmonic series $\displaystyle 1 - \frac{1}{2} + \frac{1}{3} -
    \frac{1}{4} + \cdots$ converges due to the Alternating Series Test.  Based on your answers to parts (a) and (b), what is its sum?
    \end{problem}

    \begin{solution}
    As we said in the previous part, the equation
    \[ \ln(1 + x) = \sum_{k=0}^\infty (-1)^k \frac{x^{k + 1}}{k + 1} = x
    - \frac{x^2}{2} + \frac{x^3}{3} - \cdots\] is true when $x = 1$.
    Plugging in $x = 1$ gives 
    \[ \ln 2 = \sum_{k=0}^\infty (-1)^k \frac{1}{k + 1} = 1 -
    \frac{1}{2} + \frac{1}{3} - \cdots.\]
    That is, the sum of the alternating harmonic series is exactly equal to $\boxed{\ln
    2}$.  This is something you probably wouldn't have guessed just
    from looking at partial sums of the series.
    \end{solution}
    
  \end{enumerate}

\end{enumerate}

\begin{problemonly} 
\begin{framed}
\textbf{Theorem.} If $f(x)$ has a power series representation centered
at $c$, then that power series \textit{must} be the Taylor
series of $f(x)$ centered at $c$.
\end{framed}
\end{problemonly}
\pspace{\newpage}

\begin{enumerate}[resume]
\item 
\begin{enumerate}
  \item   \note{Although they can use the Ratio Test to find the radius of convergence, I'd like them
  to see that they don't actually need to do that if they keep track
  of where the series came from.}
  
  \begin{problem}[\vfill]
  On the next PSet you will (finally) show that $\displaystyle \sin x =  x - \frac{x^3}{3!} +
  \frac{x^5}{5!} - \cdots= \sum_{k=0}^\infty
  (-1)^k \frac{x^{2k + 1}}{(2k + 1)!}$ for all $x$. Take that fact for granted now to find the Taylor series of $f(x) = x \sin (x^3)$ centered at $0$. What is its radius of convergence?
  \end{problem}

  \begin{solution}
  If we had been asked to do this before learning the theorem, we
  would have had only one option: use the fact that the Taylor series
  has the form $a_0 + a_1 x + a_2 x^2 + a_3 x^3 + \cdots$ (since we're
  centering at 0), with the coefficient $a_k$ given by $a_k =
  \frac{f^{(k)}(0)}{k!}$.  Using this fact requires us to take
  derivatives of $f(x)$ until we see a pattern, which looks pretty
  painful.

  Instead, we'll use the theorem --- it says that, if we can find a
  power series representation of $f(x)$ at $0$, then that power series
  \textit{is} the Taylor series of $f(x)$ at $0$.  So, rather than
  trying to find the Taylor series of $f(x)$ directly by taking
  derivatives, let's look for a power series representation of $f(x) =
  x \sin(x^3)$.  (Remember that this simply means we want to write an
  equation $f = $ a power series.)

  We are told that
  \[ \sin x = \displaystyle \sum_{k=0}^\infty (-1)^k \frac{x^{2k +
  1}}{(2k + 1)!} = x - \frac{x^3}{3!} + \frac{x^5}{5!} - \cdots
  \textrm{ for all } x.\] Replacing $x$ by $x^3$ everywhere gives:
  \[ \sin x^3 = \displaystyle \sum_{k=0}^\infty (-1)^k \frac{x^{6k +
  3}}{(2k + 1)!} = x^3 - \frac{x^9}{3!} + \frac{x^{15}}{5!} -
  \cdots \textrm{ for all } x.\] Now, we multiply both sides by $x$ to
  get
  \[ x \sin x^3 = \boxed{\sum_{k=0}^\infty (-1)^k \frac{x^{6k +
  4}}{(2k + 1)!}} = x^4 - \frac{x^{10}}{3!} + \frac{x^{16}}{5!}
  - \cdots \textrm{ for all } x.\]
  
  Furthermore, since the power series representation is valid for all $x$, the radius of convergence must be $\boxed{\infty}$.
  \end{solution}


  \item

  \note{The first question is mainly just to help them recall the relationship between
  Taylor polynomials and the Taylor series. Students usually have quite a bit of trouble with the last two questions. They
  remember that we find the coefficients of the Taylor series by
  taking derivatives of $f$ and doing something with $k!$, but they
  may put this information together incorrectly.  A very common
  mistake is thinking that, to find $f'''(0)$, they should look at the
  third (non-zero) term in the power series, rather than the $x^3$
  term.  Also, they sometimes remember, ``divide by $k!$'' and end up
  dividing the Taylor coefficient by $k!$ to find the derivative.
  This is again a good opportunity to remind them of the value of
  working with equations!  (If they write an equation like,
  ``coefficient of $x^k$ = $\frac{f^{(k)}(0)}{k!}$'' and manipulate
  it, they are less likely to make mistakes.)}

  \begin{problem}[\vfill]
  What is the degree 15 Taylor polynomial approximation of $f(x)$
  centered at $0$?  (Please write out all terms.)   What is $f'''(0)$?  $f^{(4)}(0)$?
  \end{problem}

  \begin{solution}
  This is the sum of the terms in the Taylor series whose degree is
  $\leq 15$, or $\boxed{x^4 - \frac{x^{10}}{3!}}$. 
  
  To answer the last two questions, let's go back to what we said in
  {{\#5}(a)}.  First, we said that
  the Taylor series of $f(x)$ was equal to $a_0 + a_1 x + a_2 x^2 +
  a_3 x^3 + \cdots$, with $a_k = \frac{f^{(k)}(0)}{k!}$.  That is,
  $a_0 = \frac{f(0)}{0!}$, $a_1 = \frac{f'(0)}{1!}$, $a_2 =
  \frac{f''(0)}{2!}$, $a_3 = \frac{f'''(0)}{3!}$, $a_4 =
  \frac{f^{(4)}(0)}{4!}$, and so on.

  At the end of {{\#5}(a)}, we
  discovered that the Taylor series of $f(x)$ centered at 0 was
  $\displaystyle \sum_{k=0}^\infty (-1)^k \frac{x^{6k + 4}}{(2k + 1)!}
  = x^4 - \frac{x^{10}}{3!} + \frac{x^{16}}{5!} - \cdots$.  If we try
  to write this in the form $a_0 + a_1 x + a_2 x^2 + a_3 x^3 +
  \cdots$, then it is equal to $0 + 0 \cdot x + 0 \cdot x^2 + 0 \cdot
  x^3 + 1 \cdot x^4 + 0 \cdot x^5 + \cdots$.  That is, if we write it
  in the form $a_0 + a_1 x + a_2 x^2 + a_3 x^3 + \cdots$, then $a_0 =
  0$, $a_1 = 0$, $a_2 = 0$, $a_3 = 0$, $a_4 = 1$, and so on.  (A
  simple way to say all of this is to say that, when we write
  something in the form $a_0 + a_1 x + a_2 x^2 + \cdots$, then $a_k$
  is the coefficient of $x^k$.)

  Now, we can simply put these two things together to find $f'''(0)$
  and $f^{(4)}(0)$.  In the first paragraph, we said $a_3 =
  \frac{f'''(0)}{3!}$ and $a_4 = \frac{f^{(4)}(0)}{4!}$.  In the
  second paragraph, we said $a_3 = 0$ and $a_4 = 1$.  So,
  $\frac{f'''(0)}{3!} = 0$ (because they are both equal to $a_3$) and
  $\frac{f^{(4)}(0)}{4!} = 1$ (because they are both equal to $a_4$).
  Multiplying the first equation by $3!$, $\boxed{f'''(0) = 0}$.
  Multiplying the second equation by $4!$, $\boxed{f^{(4)}(0) = 4!}$. 



  \textit{If you are not convinced by the above argument:} Here's
  another way to do the problem.  We showed in (a) that, in fact,
  $f(x) = x^4 - \frac{x^{10}}{3!} + \frac{x^{16}}{5!} - \cdots$ for
  all $x$.  We can differentiate this: 
  \begin{align*}
  f(x) & = x^4 - \frac{x^{10}}{3!} + \frac{x^{16}}{5!} - \cdots \\
  f'(x) & = 4x^3 - \frac{10 x^9}{3!} + \frac{16 x^{15}}{5!} -
  \cdots \\
  f''(x) & = 4 \cdot 3 x^2 - \frac{10 \cdot 9 x^8}{3!} + \frac{16
  \cdot 15 x^{14}}{5!} - \cdots \\
  f'''(x) & = 4 \cdot 3 \cdot 2 x - \frac{10 \cdot 9 \cdot 8
  x^7}{3!} + \frac{16 \cdot 15 \cdot 14 x^{13}}{5!} - \cdots \\
  f^{(4)}(x) & = 4 \cdot 3 \cdot 2 - \frac{10 \cdot 9 \cdot 8 \cdot
  7 x^6}{3!} + \frac{16 \cdot 15 \cdot 14 \cdot 13 x^{12}}{5!} -
  \cdots
  \end{align*}
  If we plug $x = 0$ in the last two equations, we get $f'''(0) = 0$
  and $f^{(4)}(0) = 4 \cdot 3 \cdot 2 = 4!$.  (This method works, but
  it would not be very fun to do if you wanted to know, say,
  $f^{(100)}(0)$.) 
  \end{solution}  
\end{enumerate}




\item 
Consider the function $f(x) = \displaystyle \frac{1}{1 - 4x^2}$.

\begin{enumerate}
  \item
  \begin{problem}[\vfill]
  Find a power series representation of $\displaystyle \int \frac{1}{1
  - 4x^2}~dx$.  For what $x$ is your representation valid?  (Give the
  largest open interval.)
  \end{problem}

  \begin{solution}
  We start with the equation
  \[ \frac{1}{1 - x} = 1 + x + x^2 + \cdots = \sum_{k=0}^\infty x^k 
  \textrm{ when } |x| < 1.\]
  Replacing $x$ by $4x^2$ everywhere in this equation,
  \[ \frac{1}{1 - 4x^2} = 1 + 4x^2 + (4x^2)^2 + \cdots =
  \sum_{k=0}^\infty (4x^2)^k \textrm{ when} \left| 4x^2 \right| < 1.\]
  Simplifying, \[ \frac{1}{1 - 4x^2} = 1 + 4x^2 + 16x^4 + \cdots =
  \sum_{k=0}^\infty 4^k x^{2k} \textrm{ when } |x| < \frac{1}{2}.\]
  Since this equation is valid for $|x| < \frac{1}{2}$, the largest
  open interval on which this representation is valid is $\left( -
  \frac{1}{2}, \frac{1}{2} \right)$.  Integrating,
  \[ \int \frac{1}{1 - 4x^2}~dx = \boxed{C + x + \frac{4}{3} x^3 +
  \frac{16}{5} x^5 + \cdots = C + \sum_{k=0}^\infty \frac{4^k}{2k +
  1} x^{2k + 1}}.\] Integrating (or differentiating) a power
  series doesn't change the largest open interval of convergence, so
  the largest open interval on which this is valid is still
  $\boxed{\left( - \frac{1}{2}, \frac{1}{2} \right)}$.  (It's
  \textit{possible} that the representation is valid at $-\frac{1}{2}$
  or $\frac{1}{2}$, but since the problem only asks for an open
  interval, we don't need to check whether the power series converges
  at $\pm \frac{1}{2}$.)
  \end{solution}

  \item 
  \begin{problem}[\vfill]
  If possible, use your answer to
part (a) to write down a series
  that converges to $\displaystyle\int_{0}^{1.5} \frac{1}{1-
  4x^2}~dx$.  If it is not possible, explain why.
  \end{problem}

  \begin{solution}
  We would like to just plug $0$ and $1.5$ in to the series we found
  in {{\#5}(a)}, but that's not valid
  because that series definitely doesn't converge when $x = 1.5$ (the
  largest open interval of convergence was $(-0.5, 0.5)$).  So, it's
  \fbox{not possible} to use the power series to evaluate
  $\displaystyle\int_{0}^{1.5} \frac{1}{1- 4x^2}~dx$.
  \end{solution}

\end{enumerate}

\end{enumerate}




\end{document}